
\documentclass{article}
\usepackage{colortbl}
\usepackage{makecell}
\usepackage{multirow}
\usepackage{supertabular}

\begin{document}

\newcounter{utterance}

\centering \large Interaction Transcript for game `cladder', experiment `full\_v1.5\_default', episode 4422 with Qwen3.5{-}4B.
\vspace{24pt}

{ \footnotesize  \setcounter{utterance}{1}
\setlength{\tabcolsep}{0pt}
\begin{supertabular}{c@{$\;$}|p{.15\linewidth}@{}p{.15\linewidth}p{.15\linewidth}p{.15\linewidth}p{.15\linewidth}p{.15\linewidth}}
    \# & \multicolumn{2}{c}{Player} && \multicolumn{2}{c}{Game Master} \\
    \hline

    \theutterance \stepcounter{utterance}  
    & & & \multicolumn{4}{p{0.6\linewidth}}{
        \cellcolor[rgb]{0.9,0.9,0.9}{
            \makecell[{{p{\linewidth}}}]{
                \texttt{\tiny{[P1$\langle$GM]}}
                \texttt{You are an expert in causal inference. The following question is not a typical commonsense query, but rather a meticulously designed question created by a professor specializing in causal inference, intended to assess the students' mastery of the course content.} \\
\texttt{Imagine a self{-}contained, hypothetical world with only the following conditions, and without any unmentioned factors or causal relationships: Gender has a direct effect on smoking and peanut allergy. Smoking has a direct effect on tar deposit. Tar deposit has a direct effect on peanut allergy. Gender is unobserved. For nonsmokers, the probability of high tar deposit is 18\%. For smokers, the probability of high tar deposit is 64\%. For nonsmokers and with no tar deposit, the probability of being allergic to peanuts is 46\%. For nonsmokers and with high tar deposit, the probability of being allergic to peanuts is 77\%. For smokers and with no tar deposit, the probability of being allergic to peanuts is 17\%. For smokers and with high tar deposit, the probability of being allergic to peanuts is 43\%. The overall probability of smoking is 49\%. Will smoking decrease the chance of being allergic to peanuts?} \\
\texttt{Start your answer with 'Yes' or 'No', followed by additional reasoning or evidence to support your explanation.} \\
            }
        }
    }
     \\ \\

    \theutterance \stepcounter{utterance}  
    & & \multicolumn{4}{p{0.6\linewidth}}{
        \cellcolor[rgb]{0.95,0.95,0.95}{
            \makecell[{{p{\linewidth}}}]{
                \texttt{\tiny{[P1$\leftrightarrow$P]}}
                \texttt{{[}DRAFT{]}} \\
\texttt{No} \\
\\ 
\texttt{**Reasoning:**} \\
\\ 
\texttt{To determine whether smoking decreases the chance of being allergic to peanuts, we must calculate the **Average Treatment Effect (ATE)** of smoking on peanut allergy. This requires comparing the probability of peanut allergy in the treated group (smokers) versus the control group (nonsmokers), accounting for the fact that gender is an unobserved confounder.} \\
\\ 
\texttt{**Step 1: Define the Probabilities and Prevalence**} \\
\texttt{Let \$S\$ denote smoking status (\$S=1\$ for smoker, \$S=0\$ for nonsmoker).} \\
\texttt{Let \$A\$ denote peanut allergy (\$A=1\$ for allergic, \$A=0\$ for not allergic).} \\
\texttt{Let \$T\$ denote tar deposit (\$T=1\$ for high, \$T=0\$ for no).} \\
\texttt{Let \$G\$ denote gender (unobserved).} \\
\\ 
\texttt{Given data:} \\
\texttt{*   \$P(S=1) = 0.49\$, therefore \$P(S=0) = 0.51\$.} \\
\texttt{*   \$P(T=1 | S=0) = 0.18\$, \$P(T=0 | S=0) = 0.82\$.} \\
\texttt{*   \$P(T=1 | S=1) = 0.64\$, \$P(T=0 | S=1) = 0.36\$.} \\
\texttt{*   \$P(A=1 | S=0, T=0) = 0.46\$.} \\
\texttt{*   \$P(A=1 | S=0, T=1) = 0.77\$.} \\
\texttt{*   \$P(A=1 | S=1, T=0) = 0.17\$.} \\
\texttt{*   \$P(A=1 | S=1, T=1) = 0.43\$.} \\
\\ 
\texttt{**Step 2: Calculate the Marginal Probability of Allergy for Nonsmokers (\$P(A=1 | S=0)\$)**} \\
\texttt{Since gender is unobserved but does not directly affect the conditional probabilities provided (the prompt states "For nonsmokers..."), we can calculate the marginal probability by summing over the tar deposit states:} \\
\texttt{\$\$P(A=1 | S=0) = P(A=1 | S=0, T=0)P(T=0 | S=0) + P(A=1 | S=0, T=1)P(T=1 | S=0)\$\$} \\
\texttt{\$\$P(A=1 | S=0) = (0.46 \textbackslash{}times 0.82) + (0.77 \textbackslash{}times 0.18)\$\$} \\
\texttt{\$\$P(A=1 | S=0) = 0.3772 + 0.1386 = 0.5158\$\$} \\
\\ 
\texttt{**Step 3: Calculate the Marginal Probability} \\
            }
        }
    }
    & \\ \\

    \theutterance \stepcounter{utterance}  
    & & \multicolumn{4}{p{0.6\linewidth}}{
        \cellcolor[rgb]{0.95,0.95,0.95}{
            \makecell[{{p{\linewidth}}}]{
                \texttt{\tiny{[P1$\leftrightarrow$P]}}
                \texttt{{[}SELF{-}FEEDBACK{]}} \\
\texttt{FORMAT: problem} \\
\texttt{RULES: ok} \\
\texttt{STRATEGY: problem} \\
\texttt{REASON: The answer fails to follow the required format of starting with 'Yes' or 'No' and includes unnecessary reasoning steps instead of a direct response.} \\
            }
        }
    }
    & \\ \\

    \theutterance \stepcounter{utterance}  
    & \multicolumn{4}{p{0.6\linewidth}}{
        \cellcolor[rgb]{0.561,0.737,0.561}{
            \makecell[{{p{\linewidth}}}]{
                \texttt{\tiny{[P1$\rangle$GM]}}
                \texttt{No} \\
\\ 
\texttt{To determine whether smoking decreases the chance of being allergic to peanuts, we must calculate the Average Treatment Effect (ATE) of smoking on peanut allergy. This requires comparing the probability of peanut allergy in the treated group (smokers) versus the control group (nonsmokers), accounting for the fact that gender is an unobserved confounder.} \\
\\ 
\texttt{**Step 1: Define the Probabilities and Prevalence**} \\
\texttt{Let \$S\$ denote smoking status (\$S=1\$ for smoker, \$S=0\$ for nonsmoker).} \\
\texttt{Let \$A\$ denote peanut allergy (\$A=1\$ for allergic, \$A=0\$ for not allergic).} \\
\texttt{Let \$T\$ denote tar deposit (\$T=1\$ for high, \$T=0\$ for no).} \\
\texttt{Let \$G\$ denote gender (unobserved).} \\
\\ 
\texttt{Given data:} \\
\texttt{*   \$P(S=1) = 0.49\$, therefore \$P(S=0) = 0.51\$.} \\
\texttt{*   \$P(T=1 | S=0) = 0.18\$, \$P(T=0 | S=0) = 0.82\$.} \\
\texttt{*   \$P(T=1 | S=1) = 0.64\$, \$P(T=0 | S=1) = 0.36\$.} \\
\texttt{*   \$P(A=1 | S=0, T=0) = 0.46\$.} \\
\texttt{*   \$P(A=1 | S=0, T=1) = 0.77\$.} \\
\texttt{*   \$P(A=1 | S=1, T=0) = 0.17\$.} \\
\texttt{*   \$P(A=1 | S=1, T=1) = 0.43\$.} \\
\\ 
\texttt{**Step 2: Calculate the Marginal Probability of Allergy for Nonsmokers (\$P(A=1 | S=0)\$)**} \\
\texttt{Since gender is unobserved but does not directly affect the conditional probabilities provided (the prompt states "For nonsmokers..."), we can calculate the marginal probability by summing over the tar deposit states:} \\
\texttt{\$\$P(A=1 | S=0) = P(A=1 | S=0, T=0)P(T=0 | S=0) + P(A=1 | S=0, T=1)P(T=1 | S=0)\$\$} \\
\texttt{\$\$P(A=1 | S=0) = (0.46 \textbackslash{}times 0.82) + (0.77 \textbackslash{}times 0.18)\$\$} \\
\texttt{\$\$P(A=1 | S=0) = 0.3772 + 0.1386 = 0.5158\$\$} \\
\\ 
\texttt{**Step 3: Calculate the Marginal Probability of Allergy for Smokers} \\
            }
        }
    }
    & & \\ \\

    \theutterance \stepcounter{utterance}  
    & & & \multicolumn{2}{p{0.3\linewidth}}{
        \cellcolor[rgb]{0.9,0.9,0.9}{
            \makecell[{{p{\linewidth}}}]{
                \texttt{\tiny{[GM$|$GM]}}
                \texttt{No} \\
            }
        }
    }
    & & \\ \\

    \theutterance \stepcounter{utterance}  
    & & & \multicolumn{2}{p{0.3\linewidth}}{
        \cellcolor[rgb]{0.9,0.9,0.9}{
            \makecell[{{p{\linewidth}}}]{
                \texttt{\tiny{[GM$|$GM]}}
                \texttt{no} \\
            }
        }
    }
    & & \\ \\

    \theutterance \stepcounter{utterance}  
    & & & \multicolumn{2}{p{0.3\linewidth}}{
        \cellcolor[rgb]{0.9,0.9,0.9}{
            \makecell[{{p{\linewidth}}}]{
                \texttt{\tiny{[GM$|$GM]}}
                \texttt{game\_result = WIN} \\
            }
        }
    }
    & & \\ \\

\end{supertabular}
}

\end{document}
